\documentclass{article}
\usepackage{PRIMEarxiv} % Core package for the PrimeAI Template
\usepackage{amsmath, amssymb, amsthm, bm, dcolumn} % Add amsthm here for the proof environment
\usepackage[numbers,sort&compress]{natbib} % Natbib for citations
\usepackage{graphicx} % For high-quality images
\usepackage{multicol} % Multi-column figures/tables
\usepackage{hyperref} % Hyperlinks
\usepackage{listings} % Code listings
\usepackage{authblk} % For structured affiliations
% Define theorem style (optional)
\theoremstyle{plain}
\newtheorem{theorem}{Theorem}
\renewcommand{\qedsymbol}{}

% Custom commands for primes and qubit encoding
\newcommand{\primeQubit}[1]{|p_{#1}\rangle}
\newcommand{\primeGate}[1]{U_{p_{#1}}}

\usepackage{wrapfig}
\usepackage[pscoord]{eso-pic}
\usepackage[fulladjust]{marginnote}
\reversemarginpar

% Typesetting improvements without footnote patching
\usepackage[protrusion=true, expansion=true, tracking=false]{microtype}
\microtypecontext{spacing=nonfrench}

% Line numbers
\usepackage[right]{lineno}

% Text layout - adjust as needed
\raggedright
\setlength{\parindent}{0.5cm}
\textwidth 5.25in 
\textheight 8.75in

% Set double spacing
\usepackage{setspace} 
\doublespacing

% Adjust width for specific content
\usepackage{changepage}

% Adjust caption style
\usepackage[aboveskip=1pt,labelfont=bf,labelsep=period,singlelinecheck=off]{caption}

% Remove brackets from references
\makeatletter
\renewcommand{\@biblabel}[1]{\quad#1.}
\makeatother

% Header, footer, and page numbers
\usepackage{lastpage,fancyhdr}
\pagestyle{fancy}
\fancyhf{}
\fancyhead[L]{Preprint - PrimeAI Enhanced Template}
\fancyfoot[C]{\scriptsize Multiplicity Theory © 2025 Ryan Van Gelder - Citizen Gardens \\ Licensed Under MIT and CC BY-NC-SA 4.0.}
\fancyfoot[R]{Page \thepage\ of \pageref{LastPage}}
\renewcommand{\footrule}{\hrule height 2pt \vspace{2mm}}

\begin{document}
\title{A Unified Spacetime Model Using Multiplicity Theory: Oscillatory Primes and Recursive Feedback}

\author{Ryan Van Gelder}
\affil{Citizen Gardens - The Foundation of Multiplicity}
\date{\today}
\maketitle

\begin{abstract}
This article presents a unified spacetime model leveraging Multiplicity Theory by integrating oscillatory states of primes, recursive feedback mechanisms, and quantum-inspired dynamics. The approach bridges classical deterministic models and quantum probabilistic frameworks, offering insights into emergent phenomena like black holes, quantum foam, and spacetime curvature. This framework, supported by prime-based encoding, tensor networks, and adaptive equations, paves the way for advanced simulations in physics, computation, and cosmology.
\end{abstract}

\section{Introduction}
Multiplicity Theory emphasizes the interconnectedness and emergent behaviors underlying the universe's fabric. This work integrates principles of prime-based oscillatory dynamics with recursive feedback to model spacetime evolution across quantum and cosmological scales. The approach seeks to unify classical and quantum paradigms, incorporating tensor networks and stochastic elements.

\subsection{Motivation and Background}
Classical physics, dominated by General Relativity, describes spacetime as a smooth manifold influenced by mass and energy. Conversely, quantum mechanics introduces granularity and probabilistic behavior. Unifying these frameworks has remained elusive. Multiplicity Theory's emphasis on eigenvalue multiplicity, recursive feedback, and holism offers a pathway to reconcile these paradigms.

\section{Mathematical Framework}

\subsection{Dynamic Multiplicity Equation}
The core of the spacetime model is expressed using a dynamic multiplicity equation:
\begin{equation}
\frac{\partial \rho_k}{\partial t} = \alpha_k(t) \rho_k + \beta_k(t) I_k + \gamma_k(t) \sum_j T_{kj} \rho_j + \lambda(t) \big(\Omega_B(\rho) + \Omega_{FS}(\rho)\big) + \xi_k(t),
\end{equation}
where:
\begin{itemize}
    \item $\rho_k$ represents the density function of the $k$-th state.
    \item $\alpha_k(t), \beta_k(t), \gamma_k(t), \lambda(t)$ are time-dependent parameters.
    \item $T_{kj}$ encodes tensor interactions between states.
    \item $\xi_k(t)$ captures stochastic noise reflecting quantum uncertainties.
    \item $\Omega_B(\rho), \Omega_{FS}(\rho)$ are geometric feedback terms, adapting to evolving spacetime states.
\end{itemize}

\subsection{Prime-Based Encoding}
Prime numbers serve as eigenvalues encoding discrete quantum states:
\begin{equation}
\phi_k = e^{2\pi i n_k / p_k} \cdot e^{i\theta_k},
\end{equation}
where $p_k$ is the $k$-th prime, and $\theta_k$ introduces phase shifts. This encoding ensures unique state representations and efficient interactions across quantum and cosmological scales.

\subsection{Tensor Networks for Multiscale Modeling}
The tensor representation of spacetime dynamics is given by:
\begin{equation}
T_{ijk} = \sum_{m,n} C_{mn} \rho_m \rho_n \otimes \phi_{ijk},
\end{equation}
where $C_{mn}$ represents coupling coefficients, and $\phi_{ijk}$ encodes prime-based interactions at higher dimensions.

\section{Applications}

\subsection{Black Hole Dynamics}
The recursive feedback mechanism is critical for modeling black hole evaporation and quantum corrections:
\begin{equation}
\Psi_{\text{BH}}(t) = \sum_{i,j} T_{ij} \Psi_i \otimes \Psi_j e^{i(\theta_i(t) + \theta_j(t))}.
\end{equation}
This equation captures quantum coherence and entanglement near event horizons, supporting simulations of Hawking radiation and the resolution of information paradoxes.

\subsection{Quantum Foam and Early Universe}
The stochastic terms $\xi_k(t)$ simulate quantum foam, where spacetime fluctuations dominate. Coupling these terms with recursive dynamics provides a framework for exploring the universe's inflationary period:
\begin{equation}
\frac{\partial \rho_k}{\partial t} = \alpha_k \rho_k + \beta_k \sum_j \xi_j(t) + \lambda \big(\Omega_B(\rho) + \Omega_{FS}(\rho)\big).
\end{equation}

\subsection{Spacetime Curvature and Tensor Dynamics}
Using the Ricci curvature tensor $R_{\mu \nu}$ and prime-based eigenvalues, spacetime can be modeled as:
\begin{equation}
R_{\mu \nu} - \frac{1}{2} g_{\mu \nu} R + \Lambda g_{\mu \nu} + \Delta_{\mu \nu} Q_{\mu \nu} = 8 \pi G \langle T_{\mu \nu} \rangle_{\text{quantum}},
\end{equation}
where $\Delta_{\mu \nu} Q_{\mu \nu}$ introduces quantum corrections influenced by prime-encoded states.

\section{Future Directions}
Future work will focus on extending this framework for high-fidelity simulations of quantum gravity, cosmological inflation, and multi-scale interactions in astrophysical phenomena. Integrating this model with real-time hypergraph dynamics and hybrid computational systems will unlock new possibilities in quantum computing and theoretical physics.

\section{Conclusion}
The integration of Multiplicity Theory into spacetime modeling offers a robust framework for unifying quantum and classical paradigms. By employing prime-based encoding, recursive feedback mechanisms, and tensor networks, this approach advances our understanding of spacetime and its emergent phenomena.


\nocite{*}
\bibliographystyle{plain}
\bibliography{references}

\end{document}